% FTR: global pointed-monoid unit, integer multiplicity, and the actual Weil receiver.
\section{The global pointed-monoid unit in the original arithmetic trace}
\label{sec:FTR-unittrace}

\paragraph{Objects and source scope.}
Let $M=\{0,\tau\}$ have absorbing element $0$ and multiplicative
unit $\tau$, so $\tau\tau=\tau$ and $0\tau=0$.
The element $\tau$ is ungraded. This is the pointed multiplicative
monoid used for the absolute base by Alain Connes and Caterina
Consani, \href{https://arxiv.org/abs/2606.06604v1}{\emph{On the
Absolute Geometry of $\operatorname{Spec}\mathbb Z$}}, original
author file \nolinkurl{FF.tex}, introduction footnote at line 145;
their monoid-to-spherical-algebra comparison appears at lines
375--405. Its unit is here named $\tau$ to match the requested
global unit. No addition on $M$ follows from its pointed
multiplication. Whenever the extra Boolean law
$\tau+\tau=\tau$ is considered below, it is stated as that
additional operation.

The actual arithmetic receivers are SLW1--29 in the supported
local Weil return chapter and SGW1--21 in the global assembly
chapter.
Their human global trace input is Alain Connes, Caterina Consani
and Matilde Marcolli,
\href{https://arxiv.org/abs/math/0703392v1}{\emph{The Weil proof
and the geometry of the adeles class space}}, original author
\nolinkurl{Yuri70v2.tex}, equation \texttt{TraceformulaH1} and
\texttt{intalphae}, lines 2014--2035.
The analysis here retains that theorem's complete test, its local
principal values, all negative as well as positive periods, and
its real and endpoint terms. It does not assign an additive trace
to a pointed monoid or an operator trace to an identity on an
infinite-dimensional space.
The real formula below is the full SGW14 comparison with Alain
Connes, \href{https://arxiv.org/abs/math/9811068v1}{\emph{Trace
formula in noncommutative geometry and the zeros of the Riemann
zeta function}}, Appendix II(34)--(37); the finite reference
principal value is the SLW18 comparison with Appendix II(9),
Lemma 2, and (22)--(29). Those existing source comparisons are
the inputs here; this derivation does not claim a new original-source
reading of that appendix or replace its principal-value constants.

\subsection{Scalar extension and expression multiplicity are exact different maps}

For a finite pointed set $k_+=\{*,1,\ldots,k\}$, the spherical
algebra $S(M)$ has pointed carrier
$M\wedge k_+=\{*,(\tau,1),\ldots,(\tau,k)\}$.
Since $\tau$ is the only nonzero element, this is the sphere
unit object on finite pointed sets. The explicit scalar map to
the integer spherical algebra is
\begin{equation}\tag{FTR1}\label{eq:FTR-sphericalunit}
 \begin{aligned}
 \eta_k:S(M)(k_+)&\longrightarrow(H\mathbb Z)(k_+)=\mathbb Z^k,\\
 *&\longmapsto(0,\ldots,0),\\
 (\tau,j)&\longmapsto(0,\ldots,\underset{j}{1},\ldots,0).
 \end{aligned}
\end{equation}
For a pointed map $\alpha:k_+\to l_+$, the target sends a vector
to the sums of its coordinates in each non-basepoint fibre of
$\alpha$. A vector in the image of \eqref{eq:FTR-sphericalunit}
has at most one nonzero coordinate, so those sums either move its
$1$ to its designated fibre or kill it when that coordinate maps
to the basepoint. This proves naturality. Pairwise multiplication
of two occupied coordinates gives $1\cdot1=1$ at the paired
coordinate, matching $\tau\tau=\tau$. Basepoints map to zero.
Thus \eqref{eq:FTR-sphericalunit} is the pointed, unital
multiplicative scalar extension, with no Boolean addition asserted.

In degree one the underlying map is
\begin{equation}\tag{FTR2}\label{eq:FTR-unitaction}
 M\longrightarrow\mathbb Z\longrightarrow\mathbb C
       \longrightarrow\operatorname{End}_{\mathbb C}(V),
 \qquad0\longmapsto0\longmapsto0\longmapsto0_V,
 \qquad\tau\longmapsto1\longmapsto1\longmapsto I_V.
\end{equation}
The last map sends a scalar $a$ to $aI_V$, and its multiplication
is composition. Hence the global scalar unit acts as identity on
every complex vector space $V$, with inverse itself.
No decomposition or grading of $V$ is required for this action.

The exact map that counts repeated occurrences is instead
\begin{equation}\tag{FTR3}\label{eq:FTR-countaction}
 \mathbb N\longrightarrow\operatorname{End}_{\mathbb C}(V),
       \qquad m\longmapsto mI_V,
 \qquad (m+n)I_V=mI_V+nI_V,\quad
       (mn)I_V=(mI_V)(nI_V).
\end{equation}
The spherical description retains the same counts: in
$(H\mathbb Z)(m_+)$ the vector $(1,\ldots,1)$ is sent by the
fold $m_+\to1_+$ to the integer $m$.
For $m\geq2$ this vector is not in the image of
$\eta_m$, whose vectors have at most one occupied coordinate.
Thus the fold that records multiplicity does not assert that an
element representing $\tau+\tau$ existed in the original
one-unit carrier with a prescribed idempotent addition.

Let $\mathbb B=\{0,\tau\}$ have Boolean addition and the same
multiplication as $M$. The presence projection is the semiring
map
\begin{equation}\tag{FTR4}\label{eq:FTR-presence}
 \pi:\mathbb N\longrightarrow\mathbb B,
       \qquad\pi(0)=0,\quad\pi(m)=\tau\ (m\geq1).
\end{equation}
For addition, $m+n$ is positive exactly when at least one
summand is positive; for multiplication, $mn$ is positive exactly
when both factors are positive. These two facts prove its laws.
If $V\ne0$, the count action \eqref{eq:FTR-countaction} does
not factor as $\mathbb N\xrightarrow\pi\mathbb B\to
\operatorname{End}(V)$: such a factorization would give
$I_V=2I_V$ from $\pi(1)=\pi(2)$, which is impossible after
subtracting $I_V$ and applying it to a nonzero vector.
The exact additive defect of the multiplicative unit action is
\begin{equation}\tag{FTR5}\label{eq:FTR-additivedefect}
 \rho(\tau+_{\mathbb B}\tau)
       -\rho(\tau)-\rho(\tau)
       =I_V-I_V-I_V=-I_V.
\end{equation}

This proves the precise two readings of the written expression
$(\tau+\tau)1$. On the free expression, evaluation in the count
receiver gives $(1+1)1=2$, and evaluation of its expanded
expression $\tau1+\tau1$ also gives $2$.
If its inner sum is first identified with the Boolean value
$\tau$, the multiplicative evaluation gives $\tau1=1$.
The expression-counting evaluation therefore does not descend to
that collapsed value. This states the exact equality that holds
and the exact factorization that fails; it does not impose
distributivity on a proposed nonadditive ambient object.

More generally, every additive map $a:\mathbb B\to A$ to an
abelian group satisfies $a(\tau)=0$, since
$a(\tau)=a(\tau+\tau)=2a(\tau)$ and cancellation holds in $A$.
In particular there is no zero- and unit-preserving additive
map from $\mathbb B$ to a nonzero ring. Conversely no unital
semiring map from a nonzero ring to $\mathbb B$ exists:
the equality $1+(-1)=0$ would imply
$\tau+a(-1)=0$, although a Boolean sum with $\tau$ is $\tau$.
The multiplicative map in \eqref{eq:FTR-unitaction} remains well defined;
the proved obstruction is specifically additive.

\subsection{Exact placement in every original support branch}

Retain the original bounded distributive lattice $L$ and
\begin{equation}\tag{FTR6}\label{eq:FTR-GL}
 \begin{gathered}
 G_L(\mathbb C)=(\{0\}\times L)
                   \cup(\mathbb C\times\{1_L\}),\\
 (a,\lambda)+(b,\mu)=(a+b,\lambda\vee\mu),\qquad
 (a,\lambda)(b,\mu)=(ab,\lambda\wedge\mu),\\
 \tau_G=(0,0_L),\qquad e=(0,1_L),\qquad
 \widehat1=(1,1_L).
 \end{gathered}
\end{equation}
The new global unit $\tau\in M$ is distinct from the historical
zero $\tau_G$ in this displayed object. There is the exact pointed
multiplicative map
\begin{equation}\tag{FTR7}\label{eq:FTR-pointedGL}
 j:M\longrightarrow G_L(\mathbb C),\qquad
       j(0)=\tau_G,\quad j(\tau)=\widehat1.
\end{equation}
Its unit action on every $G_L(V)$ is
\begin{equation}\tag{FTR8}\label{eq:FTR-fullbranchunit}
 j(\tau)(v,\lambda)=(v,1_L\wedge\lambda)=(v,\lambda),
 \qquad j(0)(v,\lambda)=(0,0_L)=\tau_V.
\end{equation}
This proves the action on every fixed label, including each
zero-amplitude point below the top. A zero scalar in the pointed
ambient semiring sends all labels to bottom, as the second formula
shows. In contrast the lift of the zero linear map retains its
label:
\begin{equation}\tag{FTR9}\label{eq:FTR-supportedzerolift}
 \widehat A(v,\lambda)=(Av,\lambda),\qquad
 \widehat{0_V}(v,\lambda)=(0,\lambda)=e(v,\lambda),\qquad
 \widehat{I_V}=\mathrm{id}_{G_L(V)}.
\end{equation}
All compositions obey $\widehat A\widehat B=\widehat{AB}$.
Thus the zero linear lift and ambient scalar zero are two
explicit maps, with the same amplitude and different label behavior.
The additive-group receiver $q_V(v,\lambda)=v$ intertwines
both with the zero map on $V$ and intertwines the new unit
with $I_V$. It does not recover discarded labels.

For integral scalar multiplicity the fixed-label action is
\begin{equation}\tag{FTR10}\label{eq:FTR-labelledcounts}
 (v,\lambda)\longmapsto(mv,\lambda),\quad m\in\mathbb N,
 \qquad
 (mv,\lambda)+(nv,\lambda)=((m+n)v,\lambda).
\end{equation}
The zero value here is the supported-zero lift in
\eqref{eq:FTR-supportedzerolift}. At top amplitude $v\ne0$,
the outputs for $m=1$ and $m=2$ are different, proving the same
failure of the presence factorization on the full carrier.
Equations \eqref{eq:FTR-fullbranchunit}--\eqref{eq:FTR-labelledcounts}
apply coordinatewise to the entire placewise carrier and to every
fixed label tuple in SLW1 and SGW1.
The unit leaves all return maps and test coefficients fixed.
No measure on the set of label tuples and no sum of one trace
for every label is introduced by this identity action.

If one sends the new $\tau$ to the old $\tau_G$ instead, its
action is the second formula of \eqref{eq:FTR-fullbranchunit},
so it is the constant-bottom map, not identity. Sending it to
$e$ gives the supported-zero map of
\eqref{eq:FTR-supportedzerolift}, also not identity on a
nonzero amplitude. These exact formulas prove which original
placement realizes a global scalar unit.

\subsection{Prime weights and Haar constants under the actual receiver}

For every prime $p$, nonzero integer $n$, and unit
$v\in\mathbb Z_p^\times$, SLW5--12 and SGW2--6 retain
\begin{equation}\tag{FTR11}\label{eq:FTR-primefull}
 \begin{gathered}
 A_{p,n,v}(x_\infty,x_p)
       =(p^{-n}x_\infty,p^{-n}v^{-1}x_p),\\
 C_{p,n,v}=
 \frac{\log p}{|1-p^{-n}|_\infty|1-p^{-n}v^{-1}|_p},\\
 w(p^n)C_{p,n,v}=\frac{\log p}{|1-p^nv|_p},\qquad
 w(y)=|y-1|,\\
 d^*u=\frac{\log p}{1-p^{-1}}\frac{du}{|u|_p},\qquad
 \operatorname{vol}_{du}(\mathbb Z_p)=1,\quad
 \operatorname{vol}_{d^*}(\mathbb Z_p^\times)=\log p.
 \end{gathered}
\end{equation}
These are literal point pullback, original product Haar measures,
and real-determinant receiving multiplier.
The scalar $\tau$ acts on its kernel as $I$; hence its flat-trace
coefficient is $C_{p,n,v}$.
Two occurrences recorded additively act as $2I$ on that kernel
and give $2C_{p,n,v}$. The weighted receiving coefficients are
respectively $\log p/|1-p^nv|_p$ and
$2\log p/|1-p^nv|_p$, which are unequal for every displayed
$p,n,v$. Thus this is an actual arithmetic witness for
\eqref{eq:FTR-additivedefect}, not only a finite-dimensional
operator example. For $n>0$ the local denominator is $1$;
for $n=-m<0$ it is $p^m$, giving the full pair
$\log p$ and $p^{-m}\log p$ before multiplication by tests.

For a locally constant compactly supported test
$b:\mathbb Q_p^\times\to\mathbb C$, the entire receiver is
\begin{equation}\tag{FTR12}\label{eq:FTR-localfull}
 \begin{aligned}
 W_p(b)={}&\sum_{n\ne0}\frac1{\log p}
  \int_{\mathbb Z_p^\times}
      w(p^n)C_{p,n,v}\,b(p^{-n}v^{-1})\,d^*v\\
 &+\int_{\mathbb Z_p^\times}
        \frac{b(v^{-1})-b(1)}{|1-v|_p}\,d^*v.
 \end{aligned}
\end{equation}
Its first sum is finite for this domain. The second numerator
vanishes near $1$, so both parts are defined with the original
principal-value subtraction. In particular every additive scalar
coefficient is received linearly, giving
\begin{equation}\tag{FTR13}\label{eq:FTR-localdefect}
 W_p(\rho(\tau+_{\mathbb B}\tau)b)
       -W_p(\rho(\tau)b)-W_p(\rho(\tau)b)=-W_p(b).
\end{equation}
For the concrete test $b=1_{p^{-1}\mathbb Z_p^\times}$, only
the $n=1$ shell contributes, and its unit test is zero. Hence
$W_p(b)=\log p\ne0$. The difference in
\eqref{eq:FTR-localdefect} is exactly $-\log p$.

The finite Haar constants use additive integration, not presence.
For $j\geq1$, the unit subgroup $K_j=1+p^j\mathbb Z_p$ has
mass
\begin{equation}\tag{FTR14}\label{eq:FTR-Haarpartition}
 \mu^*(K_j)=\frac{\log p}{1-p^{-1}}p^{-j},\qquad
 [\mathbb Z_p^\times:K_j]=(p-1)p^{j-1},\qquad
 (p-1)p^{j-1}\frac{\log p}{1-p^{-1}}p^{-j}=\log p.
\end{equation}
The coset count follows by reduction to units modulo $p^j$:
there are $p^j-p^{j-1}$ such residues. Each coset has the stated
equal mass by translation invariance. The last identity is the
integral of their disjoint indicator sum. For a $j$ with
$(p-1)p^{j-1}>1$, collapsing that sum to a single presence value
loses that displayed multiplicity. All factors must instead remain
in the real or complex additive target; a Boolean scalar value
alone cannot reconstruct them by an additive map, by the proof
following \eqref{eq:FTR-additivedefect}.

The zero-time receiving construction in SLW14--17 also retains
its full parameter. In the invariant real coordinate
$q=x_\infty/r$, the added real operator has kernel
$r^{-1}\rho(q)\sigma(q')$ relative to
$dx_\infty=r\,dq'$. Its trace parameter and supported density are
\begin{equation}\tag{FTR15}\label{eq:FTR-zerotime}
 \begin{gathered}
 c_{\rho,\sigma}=\int_{\mathbb R}\rho(q)\sigma(q)dq,\qquad
 d\Lambda_{0,v}=
  \frac{\rho(q)\sigma(q)}{|1-v^{-1}|_p}
          dq\,ds\,\delta_0(dx_p),\\
 \int d\Lambda_{0,v}
       =\frac{c_{\rho,\sigma}\log p}{|1-v|_p},\qquad
 U_p(b)=\frac1{c_{\rho,\sigma}\log p}
       \int_{\mathbb Z_p^\times}
       (b(v^{-1})-b(1))\left(\int d\Lambda_{0,v}\right)d^*v
 \quad(c_{\rho,\sigma}\ne0).
 \end{gathered}
\end{equation}
The subtraction makes the integrand vanish near $v=1$; no scalar
trace of the real identity is assigned there. For the original
pair $\rho(q)=\sigma(q)=e^{-\pi q^2}$ one has
$c_{\rho,\sigma}=1/\sqrt2$, so the numerator contains
$\log p/\sqrt2$ and the receiver contains
$\sqrt2/\log p$. They are not Boolean scalars.
The action of $\tau$ fixes the entire chosen operator and density.
The count $m$ multiplies the operator and its density by $m$;
with the fixed receiving denominator in
\eqref{eq:FTR-zerotime}, its contribution is $mU_p(b)$.
Replacing the chosen operator by $m$ copies and also changing
its receiving denominator would be a different receiving map,
not preservation of the original additive trace.

For a character test
$b(x)=h(|x|_p)\chi_p(x)$ with local unit conductor exponent $a$,
SLW20--26 prove the full formula
\begin{equation}\tag{FTR16}\label{eq:FTR-conductor}
 W_p(b)=
 \begin{cases}
 \displaystyle\log p\sum_{m\geq1}
  [z^{-m}h(p^m)+p^{-m}z^m h(p^{-m})],&a=0,\quad z=\chi_p(p),\\
 -a\log p\,h(1),&a>0.
 \end{cases}
\end{equation}
In the ramified row, the compact character annihilates every
nonzero-return unit average. Its exact remaining calculation is
\begin{equation}\tag{FTR17}\label{eq:FTR-conductorcount}
 -\log p-\sum_{j=1}^{a-1}(p^j-p^{j-1})
             \frac{\log p}{1-p^{-1}}p^{-j}
       =-a\log p.
\end{equation}
The subtraction sign, the $a$ occurrences, and the factors in
every filtration summand use addition and additive inverses in
the target. The presence projection sends every positive count
$a$ to the same value, and cannot supply the coefficient $a$.
For example at $p=2$, primitive characters of conductor $4$
and $8$ have exponents $2$ and $3$. The character modulo $4$
that sends $3$ to $-1$ is primitive. For modulus $8$, the values
$\chi(1)=\chi(7)=1$, $\chi(3)=\chi(5)=-1$ define a character
and do not factor modulo $4$, since $1$ and $5$ have different
values. Their respective terms are
$-2\log2\,h(1)$ and $-3\log2\,h(1)$.
This proves on actual compact characters the loss caused by
replacing conductor multiplicity with presence.

\subsection{The full global Weil form and its scalar operation defect}

Use exactly SGW7--14. Put
$C_{\mathbb Q}=\mathbb R_{>0}\times K$,
$K=\prod_p\mathbb Z_p^\times$, with measure $dy/y\,du$ and
$\int_Kdu=1$.
For all characters $\chi$ of a fixed finite quotient
$(\mathbb Z/N\mathbb Z)^\times$, let
\begin{equation}\tag{FTR18}\label{eq:FTR-globaltest}
 \begin{gathered}
 f(y,u)=\sum_\chi h_\chi(y)\chi(u),\qquad
 k(y,u)=\sum_\chi k_\chi(y)\chi(u),\qquad
 k^\sharp(g)=|g|^{-1}\overline{k(g^{-1})},\\
 H_\chi(y)=y^{-1}\int_0^\infty
              h_\chi(x)\overline{k_\chi(x/y)}dx,\qquad
 (f*k^\sharp)(y,u)=\sum_\chi H_\chi(y)\chi(u),\\
 M_h(s)=\int_0^\infty h(y)y^s\frac{dy}{y}.
 \end{gathered}
\end{equation}
All radial functions here are smooth and compactly supported in
$\mathbb R_{>0}$. Write $c_\chi$ for the primitive conductor,
$\chi_*$ for the primitive character, and
$\chi(-1)=(-1)^{\varepsilon_\chi}$.
The finite local terms, including the inverse local character
from the original embedding
$\iota_p(p^nv)=(p^{-n},u_p=v,u_q=p^{-n}\ (q\ne p))$, are
\begin{equation}\tag{FTR19}\label{eq:FTR-globalfinite}
 \mathcal W_{p,\chi}(H)=
 \begin{cases}
 \displaystyle\log p\sum_{n\geq1}
   [\chi_*(p)^nH(p^n)+p^{-n}\chi_*(p)^{-n}H(p^{-n})],
                                      &p\nmid c_\chi,\\
 -v_p(c_\chi)\log p\,H(1),&p\mid c_\chi.
 \end{cases}
\end{equation}
The complete real distribution is
\begin{equation}\tag{FTR20}\label{eq:FTR-archimedean}
 \begin{aligned}
 \mathcal A_\varepsilon(H)={}&(\log(2\pi)+\gamma)H(1)\\
 &+\lim_{\delta\downarrow0}\left\{
 \frac12\int_{\substack{y>0\\|1-y|\geq\delta}}
       \frac{H(y)}{|1-y|}dy+(\log\delta)H(1)\right\}\\
 &+\frac{(-1)^\varepsilon}{2}
           \int_0^\infty\frac{H(y)}{1+y}dy.
 \end{aligned}
\end{equation}
Thus the exact original global form, with differential idele
$\mathfrak a$ retained, is
\begin{equation}\tag{FTR21}\label{eq:FTR-globalfull}
 \begin{aligned}
 Q(f,k):={}&\operatorname{Tr}
       (\underline\vartheta_m(f*k^\sharp)|H^1)\\
 =\sum_\chi\bigg\{&1_{\chi=1}
       [M_{h_\chi}(0)\overline{M_{k_\chi}(1)}
        +M_{h_\chi}(1)\overline{M_{k_\chi}(0)}]\\
 &-\log|\mathfrak a|\,H_\chi(1)
       -\sum_p\mathcal W_{p,\chi}(H_\chi)
       -\mathcal A_{\varepsilon_\chi}(H_\chi)\bigg\}.
 \end{aligned}
\end{equation}
The finite-place formulas are the proved SLW receiving maps,
while the global equality is the cited human trace theorem
applied exactly as in SGW13. This retains the original two
endpoints, signs, conductor contributions, inverse-return
weights, real parity, constants and real principal value.

Let $T_{f,k}=\underline\vartheta_m(f*k^\sharp)|H^1$ be the
operator to which that theorem applies. The unit action is the
identity on its representation space, hence
$\rho(\tau)T_{f,k}=T_{f,k}$. Its trace is consequently
$Q(f,k)$, term by term in \eqref{eq:FTR-globalfull}.
Counting $m$ written copies gives $mT_{f,k}$ and trace
$mQ(f,k)$. No trace of $I_{H^1}$ on its own is used.
The exact Boolean-to-trace defect is
\begin{equation}\tag{FTR22}\label{eq:FTR-globaldefect}
 \begin{aligned}
 &\operatorname{Tr}\!\left(
 [\rho(\tau+_{\mathbb B}\tau)-\rho(\tau)-\rho(\tau)]T_{f,k}
                  \right)=-Q(f,k),\\
 &Q(\rho(\tau)f,\rho(\tau)k)=Q(f,k),\qquad
 Q(mf,nk)=m\overline n\,Q(f,k)\quad(m,n\in\mathbb C).
 \end{aligned}
\end{equation}
In particular doubling the first input doubles every term;
doubling both inputs multiplies every term by $4$.
To verify these identities without hiding a contribution,
\eqref{eq:FTR-globaltest} gives
$H_{mh,nk}=m\overline n H_{h,k}$ and
$M_{mh}(s)=mM_h(s)$. Each summand of
\eqref{eq:FTR-globalfinite}, each finite cutoff expression of
\eqref{eq:FTR-archimedean}, its limit, and the differential term
is complex-linear in $H$. The endpoint products acquire exactly
$m\overline n$. This proves the displayed factors on every
term of the original form. A particular numerical value of
$Q$ may vanish through its actual cancellations; that does not
alter the operator defect $-I$ or the proved nonzero local
witness \eqref{eq:FTR-localdefect}.

The additive operation is indispensable even in the real term:
the cutoff integral and $(\log\delta)H(1)$ form a signed sum
before the limit, and the negative-real contribution carries
the separate sign $(-1)^{\varepsilon_\chi}$ and factor $1/2$.
The powers of $\pi$ in the original archimedean source enter
here through the full constant $\log(2\pi)+\gamma$ retained
in \eqref{eq:FTR-archimedean}. An additive map out of the
Boolean value kills its unit by the earlier cancellation proof,
so it cannot carry any nonzero one of these complex-linear
coefficients. Keeping the original coefficients as external
real or complex weights still works; it is exactly the weighted
linear receiver just written, not their derivation from a
Boolean projected scalar alone.

\subsection{The intrinsic integer parity and the arithmetic scalar extension}

Retain integer parity intrinsically by the complete graph of
the original residue map:
\begin{equation}\tag{FTR23}\label{eq:FTR-intrinsicparity}
 \begin{gathered}
 \mathbb Z_{\mathrm{par}}
   =\{(n,\bar n):n\in\mathbb Z,\ \bar n=n\bmod2\},\\
 (m,\bar m)+(n,\bar n)=(m+n,\bar m+\bar n),\qquad
 (m,\bar m)(n,\bar n)=(mn,\bar m\bar n),\\
 n\longmapsto(n,\bar n),\qquad
 (n,\bar n)\longmapsto n,\qquad
 \tau\longmapsto(1,\bar1)\longmapsto I_V.
 \end{gathered}
\end{equation}
The sums and products in the second coordinates take place in
$\mathbb Z/2\mathbb Z$. The two maps between $\mathbb Z$ and
$\mathbb Z_{\mathrm{par}}$ are inverse ring isomorphisms:
the congruences $\overline{m+n}=\bar m+\bar n$ and
$\overline{mn}=\bar m\bar n$ prove preservation of both
operations, and the graph condition proves both compositions
are identity. In particular local integer $1$ has its intrinsic
odd value $\bar1$, and $1\cdot1=1$ has the same parity because
$\bar1\bar1=\bar1$. No vector-space super-grading or odd
degree-changing operator is assigned by this assertion.
The ungraded source unit $\tau$ maps to the integer unit
$(1,\bar1)$, whose full coefficient is the identity scalar.
Keeping the graph retains all of the integer's parity
information; it does not assign a parity to the source unit.

The real arithmetic parity in \eqref{eq:FTR-archimedean} is
the separate specified compact-character value at $-1$.
The scalar identity acts on every character sector and leaves
its own $\varepsilon_\chi$, conductor and phases unchanged.
It does not select the trivial compact character or any other
sector. Intrinsic integer parity alone does not determine that
character value: both the trivial character and an odd
character below use the same scalar unit $(1,\bar1)$.
For comparison, the odd primitive character modulo
$4$ is independent original arithmetic data: it has
$\chi(-1)=-1$, conductor $4$, unit term
$-2\log2\,H(1)$ at $p=2$, and at odd primes the full phases
$\chi(p)^n$ and $\chi(p)^{-n}$ in
\eqref{eq:FTR-globalfinite}. Its real last summand changes from
$+\tfrac12\int H(y)/(1+y)dy$ to the negative of that term.
These differences show exactly what arithmetic data an actual
character change would alter. Under the identity action, each
such character retains all these data. The intrinsic integer
parity graph and the character's arithmetic parity are both
kept in their specified maps and neither is substituted for
the other.

The ordinary coefficient ring generated under addition and
additive inverses by the image of $\tau$ in
\eqref{eq:FTR-sphericalunit} is precisely $\mathbb Z$, with
generator $1$. It has not acquired a second scalar generator.
At the pointed-monoid level $M$ has exactly one proper ideal,
$\{0\}$, which is prime since its complement $\{\tau\}$
is closed under multiplication. Every ring prime $\mathfrak p$
of $\mathbb Z$ pulls back under $M\to\mathbb Z$ to that same
ideal $\{0\}$: zero belongs to every ring prime and $1$
belongs to none. This gives the precise map on the degree-zero
prime data, rather than identifying the several arithmetic
primes with distinct elements of the unit monoid.

In the original support semiring the amplitude map is
$q:G_L(\mathbb Z)\to\mathbb Z$, $q(a,\lambda)=a$.
It is a surjective unital semiring map by the displayed
coordinate operations, and $q(j(\tau))=1$. Its exact prime
pullback is
\begin{equation}\tag{FTR24}\label{eq:FTR-arithmeticprimes}
 q^{-1}(\mathfrak p)
   =(\{0\}\times L)
            \cup\{(n,1_L):n\in\mathfrak p\}.
\end{equation}
It is prime because a product amplitude lies in $\mathfrak p$
exactly when one amplitude does, and it is proper because
$\widehat1$ is excluded. Conversely every prime containing
$e=(0,1_L)$ has this form: multiplying $e$ by every zero
label supplies that entire zero carrier, and its top amplitudes
form a ring ideal, with negatives supplied by multiplication
by $(-1,1_L)$. The ideal and prime tests then agree with
those in $\mathbb Z$.
The new unit action picks the already existing $\widehat1$,
which belongs to no proper ideal. It therefore creates no
additional arithmetic prime in this specified scalar target.
Other support primes of the full original lattice remain their
own existing points; the scalar map has not added or deleted
any element or ideal of that carrier.

This statement concerns exactly the scalar extension
$S(M)\to H\mathbb Z$ and its specified $G_L$ placement.
It does not compute the spectrum of a separate enlarged
ambient set with additional nonassociative or nondistributive
operations. Such an ambient completion is not needed to define
or to prove the map in \eqref{eq:FTR-unitaction}.

\subsection{The exact return to the original zeta function}

Use the original function
$\zeta(s)=\sum_{m=1}^\infty m^{-s}$, on $\Re s>1$.
The unchanged positive-period data are
\begin{equation}\tag{FTR25}\label{eq:FTR-zetareceiver}
 C_{p,n,1}=\frac{\log p}{p^n-1},\qquad
 (p^n-1)C_{p,n,1}=\log p,\qquad
 \sum_p\sum_{n\geq1}(p^n-1)C_{p,n,1}p^{-ns}
        =-\frac{\zeta'(s)}{\zeta(s)}.
\end{equation}
For completeness the original Euler product follows from unique
factorization: a finite product expands as the sum over integers
supported on its finite prime set, and absolute convergence of
$\sum m^{-s}$ passes to the limit. Its logarithmic series
$\sum_p\sum_{n\geq1}p^{-ns}/n$ converges absolutely and locally
uniformly on $\Re s>1$. On $\Re s\geq1+\eta$, $\eta>0$,
the derivative series is bounded by
\[
 \frac1{1-2^{-(1+\eta)}}
         \sum_{m=2}^\infty(\log m)m^{-(1+\eta)}<\infty,
\]
where integral comparison proves convergence of the last sum.
Termwise differentiation proves the final equality of
\eqref{eq:FTR-zetareceiver}, with every prime power retained.

Apply now the scalar extension \eqref{eq:FTR-unitaction}.
It sends each flow operator $U_{p,n}$ to
$I\,U_{p,n}=U_{p,n}$, and acts identically on its kernel,
both determinant factors, the base length $\log p$, its Haar
measure, and the receiving test. Thus every coefficient in
\eqref{eq:FTR-zetareceiver} is equal before and after that
operation. The exact defect of counting two occurrences and
then collapsing them is instead
\begin{equation}\tag{FTR26}\label{eq:FTR-zetadefect}
 \begin{aligned}
 &\sum_p\sum_{n\geq1}(p^n-1)
       [C_{p,n,1}-2C_{p,n,1}]p^{-ns}\\
 &\hspace{35mm}=-\sum_p\sum_{n\geq1}(\log p)p^{-ns}
       =\frac{\zeta'(s)}{\zeta(s)},\qquad\Re s>1.
 \end{aligned}
\end{equation}
For real $s>1$ the sum preceding the last equality is strictly
negative, so this operation defect is nonzero on the original
zeta receiver. Recording $m$ copies would similarly multiply
the original logarithmic derivative by $m$, whereas the single
global unit leaves it fixed.

One can also recover the original function from the proved
coefficient identity without an unspecified multiplicative
constant. For real $\sigma>1$, integrating the derivative series
from $\sigma$ to $+\infty$ gives
\begin{equation}\tag{FTR27}\label{eq:FTR-zetarecovery}
 \int_\sigma^\infty
      \sum_p\sum_{n\geq1}(\log p)p^{-nt}\,dt
    =\sum_p\sum_{n\geq1}\frac{p^{-n\sigma}}n
    =\log\zeta(\sigma).
\end{equation}
The integrands are nonnegative, so monotone convergence
justifies the interchange; the displayed convergent logarithmic
series gives the last equality. The original Dirichlet series
satisfies $\zeta(\sigma)\to1$ as $\sigma\to+\infty$ by
dominated convergence using any fixed exponent greater than
$1$. Therefore the action of the single scalar unit preserves
that constant as well, and its recovered function is exactly
the original $\zeta$ on this half-plane, rather than an
unspecified multiple. No extra Euler factor, rational prime,
or new zero set is introduced by this specified identity
operation. The entire Weil form still includes
\eqref{eq:FTR-globalfinite}--\eqref{eq:FTR-globalfull}; the
positive-period calculation does not remove its negative,
conductor, real, or endpoint terms.

The final operation defect is thus fully explicit at every
receiving level: $-I_V$ in \eqref{eq:FTR-additivedefect},
$-\log p$ for the actual local test following
\eqref{eq:FTR-localdefect}, $-Q(f,k)$ with every global term
retained in \eqref{eq:FTR-globaldefect}, and
$\zeta'(s)/\zeta(s)$ in \eqref{eq:FTR-zetadefect}.
The count expression with two occurrences has value $2$ in
the integer receiver; that receiver cannot factor through
the Boolean equation identifying those occurrences with one
unit. The multiplicative global-unit action and the exact
arithmetic trace remain related by the explicit scalar map,
without identifying their different additive operations.

\subsection{The exact overlap correction: truth and counting in the same receiver}

The Boolean presence value does have a constructive receiver,
with its overlap term retained. On $\mathbb B$ write $\vee$
for Boolean addition and $\wedge$ for its multiplication.
Let $\epsilon(0)=0$ and $\epsilon(\tau)=1$. For all four
pairs $b,c\in\mathbb B$, direct inspection gives
\begin{equation}\tag{FTR28}\label{eq:FTR-overlapscalar}
 \epsilon(b\vee c)=\epsilon(b)+\epsilon(c)
                       -\epsilon(b)\epsilon(c),\qquad
 \epsilon(b\wedge c)=\epsilon(b)\epsilon(c).
\end{equation}
If both values are $0$ the two sides vanish. If exactly one
is $\tau$, the first right side is $1+0-0=1$.
If both are $\tau$, it is $1+1-1=1$.
This proves the formulas exhaustively. Thus presence is received
as a valuation with an explicit overlap subtraction, not as
ordinary addition of independent occurrences.

For the actual operators $\rho(b)=\epsilon(b)I_V$, one has
\begin{equation}\tag{FTR29}\label{eq:FTR-overlapoperator}
 \begin{aligned}
 \rho(b\vee c)&=\rho(b)+\rho(c)-\rho(b\wedge c),\\
 \rho(b\wedge c)&=\rho(b)\rho(c).
 \end{aligned}
\end{equation}
These are equalities of operators on the full original receiver
space. For any operator $T$ for which the tested trace is
defined, multiply by $T$ and use its additive trace to obtain
\begin{equation}\tag{FTR30}\label{eq:FTR-overlaptrace}
 \operatorname{Tr}(\rho(b\vee c)T)
 =\operatorname{Tr}(\rho(b)T)+\operatorname{Tr}(\rho(c)T)
                  -\operatorname{Tr}(\rho(b\wedge c)T).
\end{equation}
For $b=c=\tau$, the intersection is the same unit, so the
formula is $\operatorname{Tr}(T)=\operatorname{Tr}(T)
+\operatorname{Tr}(T)-\operatorname{Tr}(T)$.
The expression count of two occurrences remains
$2\operatorname{Tr}(T)$. Equation
\eqref{eq:FTR-overlaptrace} identifies exactly the term removed
when those two occurrences designate the same presence.

Taking the local functional in \eqref{eq:FTR-localfull} and the
complete tested global trace in \eqref{eq:FTR-globalfull} gives
the constructive arithmetic identities
\begin{equation}\tag{FTR31}\label{eq:FTR-overlaparithmetic}
 \begin{aligned}
 W_p(\epsilon(b\vee c)h)
 &=W_p(\epsilon(b)h)+W_p(\epsilon(c)h)
                       -W_p(\epsilon(b\wedge c)h),\\
 Q(\rho(b\vee c)f,k)
 &=Q(\rho(b)f,k)+Q(\rho(c)f,k)
                       -Q(\rho(b\wedge c)f,k).
 \end{aligned}
\end{equation}
The first formula uses a local test $h$ in the domain of
\eqref{eq:FTR-localfull}. The second holds for the full global
tests in \eqref{eq:FTR-globaltest}. Their proofs are the actual
linearity and sesquilinearity proved above, together with
\eqref{eq:FTR-overlapscalar}; no support projection on the
arithmetic $H^1$ has been assumed.
In particular each finite prime coefficient, each conductor
summand, the full real cutoff expression and its limit, the
differential term, and both endpoints obey this overlap identity
separately. For the concrete local shell test following
\eqref{eq:FTR-localdefect}, it reads
$\log p=\log p+\log p-\log p$.
For a ramified test of exponent $a$ it reads
$-a\log p\,h(1)=-a\log p\,h(1)-a\log p\,h(1)
+a\log p\,h(1)$, preserving the original negative term and
the positive overlap correction.

For a finite family of presence values $b_1,\ldots,b_m$,
the complete formula is
\begin{equation}\tag{FTR32}\label{eq:FTR-finiteoverlap}
 \begin{aligned}
 \rho\!\left(\bigvee_{j=1}^m b_j\right)
 &=I_V-\prod_{j=1}^m(I_V-\rho(b_j))\\
 &=\sum_{\varnothing\ne A\subseteq\{1,\ldots,m\}}
       (-1)^{|A|+1}\prod_{j\in A}\rho(b_j).
 \end{aligned}
\end{equation}
All factors are commuting scalar idempotents. Their product
in the first line is $I_V$ precisely when all $b_j=0$, and
is zero otherwise, proving that line directly. Expanding its
finite product in the endomorphism ring proves the second line
with every intersection and sign retained.
Testing this finite identity against the original arithmetic
operators gives the same finite sum of full local or global
Weil expressions. When every $b_j=\tau$, every nonempty
intersection is $I_V$ and its coefficient sum is
$\sum_{r=1}^m(-1)^{r+1}\binom mr=1$, by the binomial
expansion of $(1-1)^m$. This is the exact passage from
counting $m$ occurrences to the truth of one repeated presence.

Finally the original measures themselves have this relation on
measurable finite-mass subsets, because their characteristic
functions satisfy
$1_{A\cup B}=1_A+1_B-1_{A\cap B}$ pointwise.
Integrating against the unchanged Haar measure proves
$\mu(A\cup B)=\mu(A)+\mu(B)-\mu(A\cap B)$.
For the disjoint cosets in \eqref{eq:FTR-Haarpartition}, every
overlap is empty and ordinary counting gives the full unit
mass. For a repeated identical coset, the overlap is that
entire coset and exactly one copy is subtracted. These two
calculations retain both truth and multiplicity within the
actual Haar and trace receivers. They prove a positive
comparison between the two operations without asserting a
Boolean-additive ring homomorphism or deleting the explicit
overlap term.
